\documentclass[11pt,a4paper]{article}

\usepackage[utf8]{inputenc}
\usepackage[T1]{fontenc}
\usepackage[english]{babel}
\usepackage{geometry}
\usepackage{amsmath,amssymb,amsthm,mathtools}
\usepackage{bm}
\usepackage{graphicx}
\usepackage{booktabs}
\usepackage{xcolor}
\usepackage{hyperref}
\usepackage{listings}
\usepackage{enumitem}
\usepackage{microtype}
\usepackage{fancyhdr}
\usepackage{caption}
\usepackage{subcaption}
\usepackage{float}

\geometry{margin=24mm}
\hypersetup{colorlinks=true, linkcolor=blue, urlcolor=blue, citecolor=blue}
\setlength{\headheight}{14pt}
\pagestyle{fancy}
\fancyhf{}
\lhead{Phase-Intention Hamiltonian}
\rhead{formal derivation note}
\cfoot{\thepage}

\lstset{
  basicstyle=\ttfamily\small,
  breaklines=true,
  frame=single,
  columns=fullflexible,
  keepspaces=true,
  showstringspaces=false,
  tabsize=2,
  language=Python
}

\newcommand{\dd}{\,\mathrm d}
\newcommand{\ii}{\mathrm i}
\newcommand{\Q}{\mathcal Q}
\newcommand{\Hh}{\mathcal H}
\newcommand{\Abe}{\mathcal A^{ABE}}
\newcommand{\Fbe}{\mathcal F^{BE}}
\newcommand{\Iop}{\hat{\mathcal I}}
\newcommand{\Wop}{\hat W}
\newcommand{\Jop}{\hat J}
\newcommand{\Hs}{\hat H_s}
\newcommand{\kappaI}{\kappa_I}
\newcommand{\epsEB}{\epsilon_{EB}}
\newcommand{\ThI}{\Theta_I}
\newcommand{\Lie}{\mathcal L}
\newcommand{\tr}{\operatorname{tr}}
\newcommand{\Dom}{\operatorname{Dom}}

\newtheorem{definition}{Definition}
\newtheorem{postulate}{Postulate}
\newtheorem{proposition}{Proposition}
\newtheorem{remark}{Remark}

\title{\textbf{A Phase-Intention Hamiltonian for Relational Dynamics}\\[2mm]
\large Aharonov--Bohm, Berry, Euler Closure, the Information/Intention Operator, Collatz Rhythm, and Free Twin-Prime Seeds}
\author{Adrian Lipa / USER007\\Formal working note prepared by Zosia / Codex GPT}
\date{June 19, 2026}

\begin{document}
\maketitle

\begin{abstract}
This document formulates a working Hamiltonian scaffold in which the free part of the Hamiltonian is not the usual free-particle kinetic term. Instead, the free part is taken to be a phase generator driven by an Information/Intention Operator, scaled by the constant $\ln 2/(24\pi)$, clocked by a Collatz-type discrete rhythm, and seeded by a free twin-prime pair. The horizontal or shape dynamics is then written covariantly with respect to a combined Aharonov--Bohm--Berry--Euler connection. A central Euler--Berry closure constraint is introduced, binding the phase defect $\epsilon$, the topological sector $D$, and the total holonomy. The document includes derivations, a canonical rotor embedding, a Floquet/discrete-time version, and reference simulations with source code.
\end{abstract}

\section*{Epistemic and technical status}
\begin{itemize}[leftmargin=*]
  \item \textbf{Status:} formal ansatz, derivation scaffold, and numerical consistency workbench.
  \item \textbf{Not claimed:} no experimental proof, no external NOEMA CURRENT/ACK, and no canon promotion without a separate receipt.
  \item \textbf{Main structural claim:} the phase sector is treated as the primitive free channel; the usual kinetic/shape term is horizontal dynamics relative to a phase connection.
  \item \textbf{Notation caveat:} $D$ denotes a topological closure sector: degree, winding sector, Chern/Euler sector, or an effective discrete closure index. If the project later fixes a narrower canonical meaning of $D$, the equations below should be interpreted by substituting that meaning into the same closure role.
\end{itemize}

\tableofcontents
\newpage

\section{Core thesis}

The standard perturbative decomposition
\begin{equation}
  H = H_0 + H_{\mathrm{int}}
\end{equation}
usually takes $H_0$ to be a free-particle or free-field Hamiltonian. The present relational-phase model uses a different split:
\begin{equation}
  \boxed{H = H_{\Phi,\mathcal I,C,s} + H_{\perp}.}
  \label{eq:main_split}
\end{equation}
Here $H_{\Phi,\mathcal I,C,s}$ is the free phase/intention sector, while $H_{\perp}$ is the horizontal remainder: shape, amplitude, potential, and deformation dynamics.

The seed is a free twin-prime pair
\begin{equation}
  \boxed{s=(p,p+2),\qquad p,p+2\in \mathbb P.}
\end{equation}
The Collatz/Kolac rhythm extracted from $s$ is denoted $\rho_s(k)$, where $k\in\mathbb N$ is the discrete orbital step. The Information/Intention Operator is denoted $\Iop_s(\tau,k)$.

The compact working Hamiltonian is
\begin{equation}
\boxed{
\Hs(\tau,k)
=
\frac{\hbar}{\Delta\tau_k}\rho_s(k)
\left[\frac{\ln2}{24\pi}\Wop_s+\delta\Iop_0(\tau)\right]
+
\frac12 g^{ab}\hat\Pi_a^{(s,k)}\hat\Pi_b^{(s,k)}
+V_s(q,\tau,k).}
\label{eq:compact_hamiltonian}
\end{equation}
The first term is the phase-free part. The second and third terms form the covariant horizontal dynamics.

\section{Relational configuration space and phase bundle}

Let $\Q$ be the relational configuration manifold, with local coordinates $q^a$. The phase is modeled as a $U(1)$ fiber:
\begin{equation}
  U(1)\hookrightarrow P\rightarrow \Q.
\end{equation}
Equivalently, quantum states are sections of a complex line bundle $L_s\to \Q$:
\begin{equation}
  \Psi_s(q,\chi,k)\in \Hh_s:=L^2(\Q,L_s).
\end{equation}
The phase coordinate $\chi\in U(1)$ is not treated as an isolated coordinate. It is measured covariantly through a phase connection:
\begin{equation}
  D_\tau\chi=\dot\chi+\Abe_a(q)\dot q^a.
  \label{eq:cov_phase_derivative}
\end{equation}
Thus phase evolution is not an additive decoration of the shape dynamics. Shape motion and phase motion are connected through the same bundle geometry.

\section{Aharonov--Bohm--Berry--Euler connection}

Define the total phase connection
\begin{equation}
\boxed{
  \Abe=\mathcal A_{AB}+\mathcal A_B+\mathcal A_E.}
  \label{eq:abe_connection}
\end{equation}
The components are
\begin{align}
  \mathcal A_{AB} &= \frac{q_e}{\hbar} A_i(x)\dd x^i,\\
  \mathcal A_B &= \ii\langle u(q)|\dd u(q)\rangle,\\
  \mathcal A_E &= s_E\omega_E.
\end{align}
Here $A_i\dd x^i$ is the electromagnetic vector potential, $\mathcal A_B$ is the Berry connection of the local internal/eigenstate frame $|u(q)\rangle$, and $\omega_E$ is the Euler/spin connection. In the spin-half sector,
\begin{equation}
  s_E=\frac12.
\end{equation}
The holonomy around a closed loop $\gamma$ is
\begin{equation}
  U_\gamma=\exp\left(\ii\oint_\gamma \Abe\right),
\end{equation}
and the accumulated geometric phase is
\begin{equation}
  \Phi_{ABE}[\gamma]
  =\oint_\gamma\Abe
  =\Phi_{AB}+\Phi_B+\Phi_E.
\end{equation}
The curvature of the Berry--Euler part is
\begin{equation}
  \Fbe=\dd(\mathcal A_B+\mathcal A_E)
  =\mathcal F_B+s_E\mathcal R_E,
\end{equation}
where $\mathcal F_B=\dd\mathcal A_B$ and $\mathcal R_E=\dd\omega_E$ in the Abelianized sector used here.

\section{Information constant, operator domain, and zero-level fluctuations}

The project constant is
\begin{equation}
  \boxed{\kappaI=\frac{\ln2}{24\pi}.}
  \label{eq:kappa}
\end{equation}
As a mathematical point, the constant alone is not yet a nontrivial dynamical operator. If one writes $\Iop=\kappaI\mathbf 1$, then $\Iop$ produces only a global phase in the corresponding evolution. A nontrivial generator requires an operator-valued channel, for instance a winding/intention operator $\Wop_s$:
\begin{equation}
  \boxed{\Iop_s=\kappaI\Wop_s.}
  \label{eq:intention_operator_static}
\end{equation}
To introduce dynamics from the zero level, write
\begin{equation}
  \boxed{\Iop_s(\tau)=\kappaI\Wop_s+\delta\Iop_0(\tau).}
  \label{eq:intention_operator_dynamic}
\end{equation}
The term $\delta\Iop_0(\tau)$ represents zero-level fluctuations. Formally it may be deterministic, stochastic, semiclassical, or operator-valued. For a unitary Hamiltonian sector the Hermitian part is used:
\begin{equation}
  \Iop_{H,s}(\tau)=\frac12\left(\Iop_s(\tau)+\Iop_s(\tau)^\dagger\right).
\end{equation}

\subsection{Meaning of a defined domain}

The domain of an operator is the set of states on which the operator is legally defined. It is not the numerical value $\ln2/(24\pi)$. In the present model a minimal domain is
\begin{equation}
\boxed{
\Dom(\Iop_s)=
\left\{\Psi_s\in L^2(\Q,L_s):
\Wop_s\Psi_s\in L^2(\Q,L_s),\quad
\Psi_s\text{ satisfies the holonomy boundary rule}
\right\}.}
\label{eq:domain}
\end{equation}
The holonomy boundary rule is
\begin{equation}
  \Psi_s(\gamma(1))
  =\exp\left(\ii\Phi_{ABE}[\gamma]\right)\Psi_s(\gamma(0)).
  \label{eq:holonomy_boundary}
\end{equation}
If spin-half closure is imposed, one may additionally use
\begin{equation}
  \Psi(\chi+2\pi)=-\Psi(\chi),
  \qquad
  \Psi(\chi+4\pi)=\Psi(\chi),
  \label{eq:spin_half_boundary}
\end{equation}
subject to the chosen spinorial convention.

\section{Twin-prime seed and Collatz/Kolac rhythm}

For a twin-prime seed $s=(p,p+2)$, define a scalar seed, for example,
\begin{equation}
  n_s=p(p+2).
\end{equation}
Then define the Collatz map
\begin{equation}
  C(n)=
  \begin{cases}
    n/2, & n\equiv0\pmod 2,\\
    3n+1, & n\equiv1\pmod 2.
  \end{cases}
\end{equation}
The Collatz orbit is
\begin{equation}
  n_{k+1}=C(n_k),\qquad n_0=n_s.
\end{equation}
The formalism only requires a positive rhythm functional
\begin{equation}
  \rho_s(k)=\rho[n_k,n_{k+1},\ldots]>0.
\end{equation}
A bounded reference choice used in the simulations is a parity/log-scale rhythm. Its role is not canonical; it is a numerical placeholder. The formal requirement is that $\rho_s(k)$ clocks the phase generator.

The accumulated intention/information phase over $K$ Collatz steps is
\begin{equation}
  \boxed{
  \ThI(s,K)=\sum_{k=0}^{K-1}\rho_s(k)\,\langle\Iop_s(k)\rangle.}
  \label{eq:intention_phase}
\end{equation}
In a fully operator-valued version, the corresponding phase propagator is ordered:
\begin{equation}
  U_\Phi(K,0)=
  \mathcal T\exp\left[-\ii\sum_{k=0}^{K-1}\rho_s(k)\Iop_s(k)\right].
  \label{eq:phase_propagator}
\end{equation}

\section{Euler--Berry constraint binding $\epsilon$, $D$, and phase}

The requested Euler--Berry constraint is most cleanly expressed as a closure equation. Let $\Sigma$ be a spanning surface with boundary $\partial\Sigma=\gamma$. Let
\begin{equation}
  \Phi_B[\gamma]=\oint_\gamma\mathcal A_B=\int_\Sigma\mathcal F_B,
\end{equation}
and, under the chosen Abelian spin/Euler convention,
\begin{equation}
  \Phi_E[\gamma]=\oint_\gamma s_E\omega_E=s_E\int_\Sigma\mathcal R_E.
\end{equation}
The total phase including Aharonov--Bohm and intention is
\begin{equation}
  \Phi_{\mathrm{tot}}(s,K,\gamma)
  =\Phi_{AB}[\gamma]+\Phi_B[\gamma]+\Phi_E[\gamma]+\ThI(s,K).
\end{equation}
Define the topological closure sector $D\in\mathbb Z$ and the phase defect $\epsEB\in\mathbb R$ by
\begin{equation}
\boxed{
  \Phi_{\mathrm{tot}}(s,K,\gamma)
  =2\pi\left(D+\epsEB\right).}
  \label{eq:euler_berry_constraint_phase}
\end{equation}
Equivalently,
\begin{equation}
\boxed{
  \epsEB
  =\frac{1}{2\pi}
  \left[
    \Phi_{AB}+\int_\Sigma\left(\mathcal F_B+s_E\mathcal R_E\right)+\ThI(s,K)
  \right]-D.}
  \label{eq:euler_berry_constraint_defect}
\end{equation}
Exact holonomic closure means
\begin{equation}
  \boxed{\epsEB=0.}
\end{equation}
A practical numerical or empirical tolerance version is
\begin{equation}
  \boxed{|\epsEB|\leq \epsilon_\star.}
\end{equation}

\subsection{Euler and Berry quantization interpretation}

If the Berry curvature belongs to a line bundle with first Chern number
\begin{equation}
  c_1(L_B)[\Sigma]=\frac{1}{2\pi}\int_\Sigma\mathcal F_B\in\mathbb Z,
\end{equation}
and if the Euler curvature obeys a Gauss--Bonnet-type convention
\begin{equation}
  \frac{1}{2\pi}\int_\Sigma\mathcal R_E=\chi(\Sigma),
\end{equation}
then the Berry--Euler part contributes
\begin{equation}
  \frac{1}{2\pi}\int_\Sigma(\mathcal F_B+s_E\mathcal R_E)
  =c_1(L_B)[\Sigma]+s_E\chi(\Sigma).
\end{equation}
For $s_E=1/2$, this contribution naturally carries half-integer structure whenever $\chi(\Sigma)$ is odd. This is the formal location where the Euler/Berry constraint can bind spin-half closure, the phase defect $\epsilon$, and the discrete sector $D$.

\section{Derivation I: canonical phase rotor embedding}

Consider the Lagrangian
\begin{equation}
\boxed{
  L_s
  =\frac12 g_{ab}(q)\dot q^a\dot q^b
  +\frac{I_\phi}{2}\left(D_\tau\chi\right)^2
  +J_{I,s}(\tau,k)D_\tau\chi
  -V_s(q,\tau,k),}
  \label{eq:lagrangian}
\end{equation}
where
\begin{equation}
  D_\tau\chi=\dot\chi+\Abe_a\dot q^a.
\end{equation}
The intention charge is
\begin{equation}
  J_{I,s}(\tau,k)=\hbar\rho_s(k)\Iop_s(\tau,k),
\end{equation}
or its expectation/scalar reduction in a semiclassical sector.

The canonical momentum conjugate to $\chi$ is
\begin{align}
  J
  &=\frac{\partial L_s}{\partial\dot\chi}\\
  &=I_\phi D_\tau\chi+J_{I,s}.
  \label{eq:J_momentum}
\end{align}
Thus
\begin{equation}
  D_\tau\chi=\frac{J-J_{I,s}}{I_\phi}.
  \label{eq:Dchi_solved}
\end{equation}
The momentum conjugate to $q^a$ is
\begin{align}
  p_a
  &=\frac{\partial L_s}{\partial\dot q^a}\\
  &=g_{ab}\dot q^b+\frac{\partial}{\partial\dot q^a}\left[\frac{I_\phi}{2}(D_\tau\chi)^2+J_{I,s}D_\tau\chi\right]\\
  &=g_{ab}\dot q^b+\bigl(I_\phi D_\tau\chi+J_{I,s}\bigr)\Abe_a\\
  &=g_{ab}\dot q^b+J\Abe_a.
  \label{eq:p_momentum}
\end{align}
Therefore the horizontal/covariant momentum is
\begin{equation}
  \Pi_a=p_a-J\Abe_a=g_{ab}\dot q^b,
  \label{eq:cov_momentum}
\end{equation}
and hence
\begin{equation}
  \dot q^a=g^{ab}\Pi_b.
\end{equation}

The Hamiltonian is
\begin{align}
  H_s
  &=p_a\dot q^a+J\dot\chi-L_s.
\end{align}
Using $\dot\chi=D_\tau\chi-\Abe_a\dot q^a$ and $p_a\dot q^a-J\Abe_a\dot q^a=\Pi_a\dot q^a$, one obtains
\begin{align}
  H_s
  &=\Pi_a\dot q^a+J D_\tau\chi
  -\frac12 g_{ab}\dot q^a\dot q^b
  -\frac{I_\phi}{2}(D_\tau\chi)^2
  -J_{I,s}D_\tau\chi+V_s\\
  &=\frac12 g^{ab}\Pi_a\Pi_b
  +\left(J-J_{I,s}\right)D_\tau\chi
  -\frac{I_\phi}{2}(D_\tau\chi)^2+V_s.
\end{align}
Substituting Eq.~\eqref{eq:Dchi_solved} gives
\begin{equation}
\boxed{
  H_s
  =\frac{(J-J_{I,s})^2}{2I_\phi}
  +\frac12 g^{ab}(p_a-J\Abe_a)(p_b-J\Abe_b)
  +V_s(q,\tau,k).}
  \label{eq:canonical_hamiltonian}
\end{equation}
This is the canonical rotor version of the model.

Its phase-free part is
\begin{equation}
  \boxed{H_{\Phi}^{\mathrm{rotor}}=\frac{(J-J_{I,s})^2}{2I_\phi}.}
\end{equation}
Its horizontal part is
\begin{equation}
  \boxed{H_\perp=\frac12g^{ab}\Pi_a\Pi_b+V_s.}
\end{equation}

\section{Derivation II: discrete Floquet phase generator}

Because the Collatz/Kolac rhythm is naturally discrete, a Floquet formulation is often cleaner. Define the phase propagator per step by
\begin{equation}
\boxed{
  U_\Phi(k+1,k)
  =\exp\left[-\ii\rho_s(k)\Iop_s(k)\right].}
  \label{eq:floquet_step}
\end{equation}
If
\begin{equation}
  U_\Phi(k+1,k)=\exp\left[-\frac{\ii}{\hbar}\Delta\tau_k H_\Phi(k)\right],
\end{equation}
then the effective phase Hamiltonian is
\begin{equation}
\boxed{
  H_{\Phi,\mathcal I,C,s}(k)
  =\frac{\hbar}{\Delta\tau_k}\rho_s(k)\Iop_s(k).}
  \label{eq:floquet_phase_hamiltonian}
\end{equation}
Substituting Eq.~\eqref{eq:intention_operator_dynamic},
\begin{equation}
\boxed{
  H_{\Phi,\mathcal I,C,s}(\tau,k)
  =\frac{\hbar}{\Delta\tau_k}\rho_s(k)
  \left[\kappaI\Wop_s+\delta\Iop_0(\tau)\right].}
  \label{eq:floquet_phase_hamiltonian_dynamic}
\end{equation}
This is the form used in Eq.~\eqref{eq:compact_hamiltonian}.

\section{Covariant quantum Hamiltonian}

Quantize the phase charge by
\begin{equation}
  \Jop=-\ii\hbar\frac{\partial}{\partial\chi}.
\end{equation}
The horizontal momentum becomes
\begin{equation}
\boxed{
  \hat\Pi_a^{(s,k)}
  =-\ii\hbar\nabla_a
  -\hbar\alpha_s(k)\Abe_a
  -\lambda_s(k)\mathcal I_a(q).}
  \label{eq:quantum_cov_momentum}
\end{equation}
Here $\alpha_s(k)$ is the effective phase charge/winding sector and $\mathcal I_a$ is the geometric trace of the intention channel on $\Q$. In a fixed $\Jop$ eigen-sector one may identify $\alpha_s(k)$ with the corresponding phase charge.

The full discrete-time Hamiltonian is
\begin{equation}
\boxed{
\Hs(\tau,k)=
\frac{\hbar}{\Delta\tau_k}\rho_s(k)
\left[\kappaI\Wop_s+\delta\Iop_0(\tau)\right]
+
\frac12 g^{ab}\hat\Pi_a^{(s,k)}\hat\Pi_b^{(s,k)}
+V_s(q,\tau,k).}
\end{equation}
The same structure in canonical rotor embedding is
\begin{equation}
\boxed{
\Hs^{\mathrm{rotor}}(\tau,k)=
\frac{(\Jop-\hbar\rho_s(k)\Iop_s(\tau,k))^2}{2I_\phi}
+
\frac12 g^{ab}(\hat p_a-\Jop\Abe_a)(\hat p_b-\Jop\Abe_b)
+V_s.}
\end{equation}

\section{Consistency conditions}

For the scaffold to be mathematically meaningful, the following conditions are required.

\begin{enumerate}[leftmargin=*]
  \item \textbf{Domain condition:} $\Dom(\Iop_s)$ must be fixed, at least by Eq.~\eqref{eq:domain}.
  \item \textbf{Hermiticity condition:} the Hamiltonian part of $\Iop_s$ must be self-adjoint or replaced by its self-adjoint part.
  \item \textbf{Rhythm condition:} $\rho_s(k)>0$ must be defined for every step in the selected Collatz orbit.
  \item \textbf{Closure condition:} the Euler--Berry defect must obey Eq.~\eqref{eq:euler_berry_constraint_defect}; exact closure means $\epsEB=0$.
  \item \textbf{Bundle condition:} the wavefunction must obey the ABE holonomy rule Eq.~\eqref{eq:holonomy_boundary}.
  \item \textbf{No false energy claim:} AB, Berry, and Euler phases are connections/holonomies. They become Hamiltonian terms only after a generator, rotor inertia, or Floquet time step is specified.
\end{enumerate}

\section{Reference simulations}

The package contains three reference simulations. They are internal consistency checks, not physical evidence.

\subsection{Collatz phase accumulation}

The script \texttt{scripts/collatz\_phase\_sim.py} uses the twin-prime seed $(17,19)$, scalar seed $n_s=323$, and $80$ Collatz steps. It computes a bounded Collatz rhythm $\rho_s(k)$, a scalar proxy for $W_s(k)$, and the accumulated intention phase. The resulting closure values are read from \texttt{results/collatz\_phase\_result.json}.

\begin{figure}[H]
\centering
\includegraphics[width=0.88\linewidth]{figures/collatz_orbit.png}
\caption{Collatz orbit from the twin-prime scalar seed $17\cdot19=323$.}
\end{figure}

\begin{figure}[H]
\centering
\includegraphics[width=0.88\linewidth]{figures/collatz_rhythm.png}
\caption{Reference bounded Collatz rhythm $\rho_s(k)$.}
\end{figure}

\begin{figure}[H]
\centering
\includegraphics[width=0.88\linewidth]{figures/intention_phase_accumulation.png}
\caption{Accumulated scalar intention/information phase $\Theta_I$.}
\end{figure}

For the default run,
\begin{align}
  \kappa_I&=0.009193150006360484,\\
  \Theta_I&=0.3033889217620433,\\
  \Phi_{\mathrm{tot}}&=4.484981575351837,\\
  D&=1,\\
  \epsilon&=-0.2861930126066793.
\end{align}
The unitary modulus error for $e^{\ii\Phi_{\mathrm{tot}}}$ is zero to displayed precision.

\subsection{Euler--Berry closure scan}

The script \texttt{scripts/euler\_berry\_constraint\_scan.py} scans Berry phase values and Euler characteristics $\chi\in\{-2,-1,0,1,2\}$. It computes
\begin{equation}
  \epsilon=\frac{\Phi_{AB}+\Phi_B+s_E2\pi\chi+\Theta_I}{2\pi}-D.
\end{equation}

\begin{figure}[H]
\centering
\includegraphics[width=0.88\linewidth]{figures/euler_berry_constraint_scan.png}
\caption{Absolute Euler--Berry closure defect across a small phase/topology grid.}
\end{figure}

For the default scan, the closest closure point is approximately
\begin{equation}
  \chi=-2,\qquad \Phi_B=-0.7330382858376185,
  \qquad D=-1,
  \qquad \epsilon\approx -4.8355820958301265\times10^{-4}.
\end{equation}

\subsection{Finite matrix Hamiltonian check}

The script \texttt{scripts/hamiltonian\_matrix\_demo.py} constructs a finite-dimensional Hermitian matrix model containing a phase term and a nearest-neighbor holonomy coupling. It checks Hermiticity, unitarity of the generated propagator, norm conservation, and energy expectation under a time-independent reference Hamiltonian.

\begin{figure}[H]
\centering
\includegraphics[width=0.88\linewidth]{figures/unitary_norm_check.png}
\caption{Norm preservation under the finite-dimensional unitary propagator.}
\end{figure}

\begin{figure}[H]
\centering
\includegraphics[width=0.88\linewidth]{figures/energy_expectation_demo.png}
\caption{Energy expectation for the time-independent finite matrix reference model.}
\end{figure}

Default numerical checks:
\begin{align}
  \|H-H^\dagger\|&=0,\\
  \|U^\dagger U-I\|&\approx6.2610\times10^{-15},\\
  \max_k\left|\|\psi_k\|-1\right|&\approx6.9944\times10^{-15}.
\end{align}
These values indicate numerical Hermiticity/unitarity in the reference finite model.

\section{Interpretation}

The formal pipeline is
\begin{equation}
\boxed{
  (p,p+2)
  \longrightarrow
  n_s
  \longrightarrow
  \{n_k\}_{k\ge0}
  \longrightarrow
  \rho_s(k)
  \longrightarrow
  \Iop_s(k)
  \longrightarrow
  \Phi
  \longrightarrow
  \text{ABE holonomic relational dynamics}.}
\end{equation}
The free part is therefore not a free-particle Hamiltonian. It is the free phase/intention generator:
\begin{equation}
\boxed{
  H_{\mathrm{free}}
  =H_{\Phi,\mathcal I,C,s}.}
\end{equation}
The ordinary-looking kinetic term belongs to the horizontal geometry and is already covariantly deformed by Aharonov--Bohm, Berry, Euler, and intention-trace connections.

\section{Minimal next formal tasks}

\begin{enumerate}[leftmargin=*]
  \item Fix the canonical project definition of $D$.
  \item Choose whether the primary Hamiltonian is the rotor embedding Eq.~\eqref{eq:canonical_hamiltonian} or the Floquet embedding Eq.~\eqref{eq:floquet_phase_hamiltonian_dynamic}.
  \item Replace scalar $W_s(k)$ in simulations with an explicit finite operator $\Wop_s$.
  \item Define the fluctuation law for $\delta\Iop_0(\tau)$: deterministic, stochastic, semiclassical, or operator-valued.
  \item Define a project-canonical Collatz rhythm $\rho_s(k)$.
  \item Add automated tests for Hermiticity, domain admissibility, closure defect, and unitary propagation.
\end{enumerate}

\appendix
\section{Source code: Collatz phase simulation}
\lstinputlisting[caption={collatz\_phase\_sim.py}]{scripts/collatz_phase_sim.py}

\section{Source code: Euler--Berry constraint scan}
\lstinputlisting[caption={euler\_berry\_constraint\_scan.py}]{scripts/euler_berry_constraint_scan.py}

\section{Source code: finite Hamiltonian matrix demo}
\lstinputlisting[caption={hamiltonian\_matrix\_demo.py}]{scripts/hamiltonian_matrix_demo.py}

\section{Source code: run all}
\lstinputlisting[caption={run\_all.py}]{scripts/run_all.py}

\end{document}
